% !TeX root = ../papers/nc-p3.tex
\section{滑らかな整五次点グラフの二同型類}
\label{sec:elliptic-quintic-regularity}
\begin{theorem}\label{thm:elliptic-quintic-regularity}
滑らかな整五次曲線とそれに交わらない孤立残余を点グラフとし、四辺の左右射影ファイバーが
空か一点となる標準14641型の添字保持同型類は、前定理の二つで尽きる。
どちらも孤立残余は空で、点グラフは被約な楕円五次曲線である。
両者とも最小の裸周期は1で、標準四移動は点曲線上で恒等である。
\end{theorem}
\begin{proof}[隣接点グラフのスキーム]
$L,M$ を一辺の二偏極、$V,W$ をその切断系とする。隣接差 $a,b$ は五捩れではないので $L\ne M$。
二つの曲線には共通の三割線因子がない。前節の三つ組のうち差が $k$ の部分対は、
重複を含めても一意であり、同じ部分対に対する第三点は $\ker(2-\epsilon)$ 外の平行移動で変わる。

$H^0(LM)^*$ の非零な $\lambda$ が $V\otimes W$ 上で階数1を与えるとする。
完全系での双線型形式は階数3以下で、命題\ref{prop:quintic-mixed-kernel}の拡大を用いると
$F L^{-1}$ は二つ以上の独立切断を持つ。行列式 $M/L$ が非自明な次数零の束なので、
これらの切断は全て一本の飽和直線束に入り、その次数は2以上となる。
従って $\lambda$ は長さ3以下の因子 $Z$ 上の汎関数である。$Z$ を最小に取れば、
その上の乗法対合は非退化となる。実際、核イデアルが非零ならそのイデアルで割ったより短い部分を通る。
長さ2までは両空間の評価が全射で、その制限対合の階数は長さに等しい。
長さ3なら両像は次元2以上であり、一方でも全射なら階数2以上である。
両像が次元2なら共通三割線ができてしまう。よって $Z$ は一点であり、
点対は曲線のグラフ上にある。逆に曲線上の全点は関係を消すので、これで台が決まる。

曲線上の $p$ で $V(-p), W(-p)$ と書いて、それぞれ $p$ で消える三切断を基点を除いて見る。
両者は $L(-p),M(-p)$ の基点自由な超平面である。その積は $H^0(LM(-2p))$ を張る。
実際、この積を消す汎関数の全双線型階数は2以下で、次数4の二偏極について上と同じ拡大を使えば
二つ以上の切断と行列式 $M/L\ne1$ により、汎関数は長さ2以下の部分に支持される。
基点自由性で長さ1は不可能。長さ2で一方の評価が全射なら他方の単元切断との積で全切断を得る。
両方とも全射でなければこの長さ2の因子に $p$ を加えたものが共通三割線となる。
階数1行列でのSegre錐の接条件は、左右の消滅超平面の積上で微分が零となることである。
従ってグラフのアフィン錐の接空間は $H^0(LM(-2p))$ の零化空間で次元2、
射影接空間は次元1である。全閉点が滑らかとなりスキームとして曲線のグラフに等しい。

\smallskip
\noindent\textit{有限Frobenius性と中心切断。}
以下は $\epsilon=\mathrm i$ で示す。共役により逆符号も従う。
$D=\operatorname{diag}(1,-1,\mathrm i,-\mathrm i)$ は曲線上で $[\mathrm i]$ に対応し、
$(D\otimes D)R(a)=R(b)$、$(D^2\otimes D^2)R(a)=R(a)$ である。
位置 $i$ の基底を $D^i$ で変えれば全辺の関係は $U=(1\otimes D)R(a)$ に一致する。
よって通常二次代数 $T(V)/(U)$ の添字展開として裸周期1を得る。

\path{computations/verify_elliptic_quintic_as.py}は $\mathbb Q(\mathrm i)$ 上で次を確認する。
三文字、四文字の隣接関係の交叉次元は4,1。四次テンソルの三つの切断の階数は4,6,4で、
循環置換はこのテンソルを固定する。これは双対の全積対合の完全性を与える。
五次交叉は零である。両端支持が4の四次テンソルに対し、その右に一文字を付けた空間と
左に一文字を付けた空間が非零に交われば、最初の一文字の支持が同時に4と1になってしまう。
これで標準有限Frobenius性が従い、循環恒等式から標準四移動も射影的に恒等である。

通常代数の四次部分で、全一次文字 $x$ に対し $u x=\lambda x u$ となる空間を
$\lambda=1,-1,\mathrm i,-\mathrm i$ の順に求めると、次元は $2,1,1,1$。
計算は単項式を長さ5まで逐次商にし、五次でこの線型方程式を解く。
各解 $u_j$（$1\le j\le5$、基底選択は同コード）に対し、頂点 $m$ で
$\lambda_j^m u_j$ とする族は折り畳み環で中心となる。
これらを $j^g$（$0\le g\le3$）で線型結合した四つを $z_g$ とする。

\smallskip
\noindent\textit{全次数の認証。}
これらの式は181で良い還元を持つ。標数181において $z_g$ による商は四始点全てで
次数0から13まで
\[
 (1,4,10,20,31,40,44,40,31,20,10,4,1,0)
\]
となる。これは同コードが各長さで、直前までの商に最後の二次関係と四つの四次中心関係を挿入して求める。

標数181での通常の道空間と左右線型複体は
\path{computations/verify_elliptic_quintic_koszul.py}で別に計算する。
交代二捩れの元の基底に戻し、位数5の指標で各行列を分ける。始点0について、長さ $n-2$ の道の
末尾に六関係を挿入し、長さ $n-1$ の道と最後の文字の積を商にする。
すべての $n\le16$ で次元は $h_n=\binom{n+3}{3}$。
線型複体の右端二箇所の完全性は一次生成・二次表示から従う。左端の微分については、
四次テンソルの一文字縮約が独立なので、接頭部の道への作用の核は四つの末尾文字の共通零化空間である。
各次数 $n\le16$ の該当する道でこの共通零化空間は零である。
中央の三次交叉から二次交叉への微分は、三次テンソルを一文字と二次関係に分けて、
その一文字を接頭部に掛けた行列である。像階数は $4h_{n-3}-h_{n-4}$
（負次数は零）となる。従ってEuler式により有限範囲で完全である。
始点を一つ進めることは $[\mathrm i]$ による基底変更であり、両隣接差は二捩れゆえ
$R(a),R(b)$ はそれぞれ転置で不変である。逆順の交代窓も同じなので反対側と全始点にこの計算が及ぶ。
還元後の三つの切断の階数と四次交叉の次元も同コードで確認する。

定理\ref{thm:central-window-cutoff}の十分条件が標数181で成立する。その証明は体の標数を用いない。
標数零への適用では、四中心式の存在を特殊化から推測する必要はない。上で標数零の式を先に解いている。
同定理の開性の証明をこの離散付値環上で用いる。具体的には、次数16まで、有限Frobenius枠の一定階数と
線型複体の単元小行列から、次の道空間が自由であることと完全性を次数帰納で持ち上げる。
中心Koszul複体は特殊ファイバーで正の位置が完全であり、その同じ階数小行列により
中心商の必要な次元も次数13まで持ち上がる。よって同定理を標数零で適用でき、
二候補はともに全次数でKoszul左右AS正則である。先の点グラフの照合と二類の非同型性を併せれば主張を得る。
\end{proof}
